\documentclass[11pt,a4paper]{article}
% Compile with XeLaTeX (fontspec and xeCJK).
\usepackage{fontspec}
\setmainfont{LibertinusSerif}[Extension=.otf,UprightFont=*-Regular,BoldFont=*-Bold,ItalicFont=*-Italic,BoldItalicFont=*-BoldItalic]
\setsansfont{LibertinusSans}[Extension=.otf,UprightFont=*-Regular,BoldFont=*-Bold,ItalicFont=*-Italic]
\usepackage{xeCJK}
\setCJKmainfont{Hiragino Mincho ProN}
\xeCJKDeclareCharClass{Default}{"00B0, "00B7, "2013, "2014, "2018 -> "201F, "2026}
\usepackage[french]{babel}
\usepackage[a4paper,margin=2.5cm,top=2.3cm]{geometry}
\usepackage{graphicx}
\usepackage{xcolor}
\usepackage{titlesec}
\usepackage{fancyhdr}
\usepackage{enumitem}
\usepackage{xurl}
\usepackage[hidelinks]{hyperref}
\definecolor{accent}{RGB}{31,73,125}
\definecolor{muted}{RGB}{93,104,119}
\hypersetup{colorlinks=true,urlcolor=accent,linkcolor=accent,pdfauthor={Romain Fernex},pdftitle={Un OTAN asiatique est-il possible ? Le Dialogue quadrilatéral pour la sécurité et les enjeux de la coopération stratégique en Indopacifique}}
\urlstyle{same}
\setlist[itemize]{topsep=.25em,itemsep=.15em,parsep=0pt,leftmargin=1.4em}
\setcounter{secnumdepth}{0}
\titleformat{\section}{\large\bfseries\color{accent}}{}{0em}{}
\titlespacing*{\section}{0pt}{1.6em}{.6em}
\setlength{\parindent}{1.2em}
\setlength{\parskip}{0.35em}
\pagestyle{fancy}
\fancyhf{}
\fancyhead[L]{\small\color{muted}Romain Fernex}
\fancyhead[R]{\small\color{muted}\itshape Un OTAN asiatique est-il possible ?}
\fancyfoot[C]{\small\color{muted}\thepage}
\renewcommand{\headrulewidth}{0pt}
\newcommand{\partbreak}{\par\begin{center}\color{muted}$\ast$\quad$\ast$\quad$\ast$\end{center}\par}


\begin{document}
\thispagestyle{empty}
\noindent{\small\sffamily\color{accent}\MakeUppercase{Série d’articles · Sécurité en Indopacifique}}\par\medskip
\noindent{\LARGE\bfseries Un OTAN asiatique est-il possible ?\par}
\smallskip\noindent{\large\itshape Le Dialogue quadrilatéral pour la sécurité et les enjeux de la coopération stratégique en Indopacifique\par}
\medskip
\noindent{\color{accent}\rule{\linewidth}{0.8pt}}\par\smallskip
\noindent{\small\color{muted}Romain Fernex · 2022 · Série d’articles en trois parties, en français}\par\bigskip

\section{Partie 1/3}

Le 29 juin dernier, à Madrid, s'est tenu le premier sommet de l'OTAN auquel ont été convié ses puissances partenaires dans l'espace Indo Pacifique. Espace stratégique d'importance capitale pour les Etats-Unis qui le définissent comme s'étendant de la côte Ouest des Etats-Unis à l'Océan Indien\textbf{{[}1{]}}, l'espace Indo Pacifique représente plus des 2/3 de l'économie mondiale et englobe plus de 50\% de la population mondiale. Son importance se fait d'autant plus grande que les tensions en Europe ont accentué le déplacement du «~centre du monde~» vers cette zone bien que celui-ci fut déjà bien entamé depuis 1985,. Ainsi, à l'occasion de ce sommet, le secrétaire général de l'OTAN, Jens Stoltenberg, a évoqué le défi posé par la guerre en Ukraine et ses effets déstabilisateurs mais il en a aussi profité pour mettre l'accent sur les tensions toujours plus préoccupantes dans le Pacifique Sud en soulignant, de manière plus ou moins voilée, le risque que constitue la montée de la Chine pour la sécurité internationale dans cette zone. Si cette position n'a rien de surprenant historiquement, étant donné qu'elle reprend largement la vision stratégique défendue par les Etats-Unis depuis plusieurs décennies, le fait d'aborder cette problématique hors du cadre traditionnel des alliances bilatérales et trilatérales qui dominent encore les relations en Asie Pacifique n'est pas anodin. De fait, pour les Etats-Unis, comme pour ses alliés, la question du „Containment „ \textbf{{[}2{]}} de la Chine passe habituellement par des accords tel que le traité de coopération mutuelle et de sécurité entre les Etats-Unis et le Japon adopté le 16 mai 1960 et qui demeure un des piliers de la présence américaine en Indo Pacifique. Cependant, Il ne s'agit pas pour autant pour ces puissances d'intégrer l'OTAN comme l'a rappelé John Kirby\textbf{{[}3{]},} coordinateur des communications stratégiques au conseil de sécurité de la maison Blanche. Dès lors, on peut se demander si, à défaut d'une extension de l'OTAN, une structure d'alliance similaire est tout de même possible en Asie en dehors du domaine strictement économique.

A ce titre, certains analystes politiques soulignent l'importance de groupements informels déjà existants et qui pourraient potentiellement, notamment dans le contexte géopolitique actuel, aboutir à la création de réelles alliances dont le rôle et la nature seraient proche de l'OTAN. Parmi ces groupements, l'un apparait plus prometteur que les autres, le Quadrilateral Security Dialogue (ou Quad). Cette initiative lancée en 2007 et ravivée en 2017, réunissant le Japon, l'Australie, L'inde et les Etats-Unis a également fait les gros titres récemment avec son quatrième sommet qui s'est tenu le 24 mai 2022 à Tokyo, soit à peine plus d'un an après le sommet précédent en mars 2021. L'accélération que connait le Quad a ainsi rapidement été sujette à critique, notamment de la part de la Chine qui l'a qualifié de «~clique qui pourrait commencer une nouvelle guerre froide~», pour reprendre les mots du ministre des affaires étrangères chinois Wang Yi en mars 2021. Pourtant le Quad dans sa forme actuelle ne possède pas de dimension militaire et son volet sécuritaire reste relativement anecdotique étant donné les différences d'opinions entre les membres du Quad, notamment entre l'Inde et les Etats-Unis, sur ce qu'il doit devenir.

Alors, qu'en est-il réellement~? Qu'est-ce que le Quad à l'heure actuelle~? Quelles sont les dynamiques qui l'animent et qui s'y font face~? et y-a-t 'il lieu d'y voir, comme l'entend la RPC et certains de ses alliés, un futur OTAN de l'Indo Pacifique~? C'est à ces questions que nous tenterons de répondre au cours de cette série de 3 articles. Le premier article traitera ainsi des origines du Quad et de sa forme actuelle. Puis dans un second article nous verrons quels sont les dynamiques qui s'opposent au sein même du Quad et en quoi la crise peut aujourd'hui nous amener à repenser son rôle. Enfin nous interrogerons la faisabilité ainsi que la légitimité de faire du Quad un~«~OTAN de l'Asie~» en présentant différents scenarii d'évolution pour celui-ci.

Tout d'abord, avant de se pencher sur la création du Quad, il s'agit de souligner que celui-ci est loin d'être la première initiative multilatérale en Indo Pacifique. Dans le domaine économique on peut citer l'ASEAN, fondée le 8 aout 1967 par la déclaration de Bangkok, qui regroupe déjà une dizaine de membres, et son regional forum qui a lieu chaque année depuis 1994 et convie 27 membres dont la majorité sont extérieurs à l'ASEAN\textbf{{[}4{]}~}comme le Japon, la Corée du Sud, la Chine, L'Inde et les Etats-Unis. Pour ce qui est du domaine stratégique, qui est celui qui nous intéresse ici, on peut également évoquer l'OTASE, fondée le 8 septembre 1954 en pleine guerre froide par les Etats-Unis, rejoint par 7 autres états, et dissoute le 30 juin 1977, qui avait vocation à être un véritable OTAN asiatique. Mais si l'OTASE a disparu avec le réchauffement des relations entre pays asiatiques en dépit de la guerre froide, le retour des tensions pose inévitablement la question du retour de formes de coopérations similaires.

L'idée Quad n'est cependant pas, à l'origine, apparue comme une réponse à ces nouvelles tensions : En effet, les bases opérationnelles du Quad ont été posées avec la campagne d'aide humanitaire menée conjointement par les futurs membres du Quad lors du boxing day tsunami qui a touché les côtes des pays d'Asie pacifique en 2004. Le but était pour ces 4 pays de coordonner leurs efforts pour mener au mieux cette opération et cela renvoie à un objectif important du Quad, y compris dans sa forme actuelle, qui est de contribuer au développement des pays du Pacifique Sud en apportant une alternative plus multilatérale à l'aide majoritairement bilatérale offerte par la Chine. Le fondement idéologique du Quad est, lui, à trouver dans un discours, connu sur le nom de «~confluence of the two seas~speech», prononcé par Shinzo Abe en 2006, alors qu'il était candidat aux élections pour le poste de premier ministre au Japon. Dans ce discours, il appel à la constitution d'un «~réseau d'états unis autour des principes de liberté et de primauté de la loi (la loi internationale ici) qui s'étendrait d'un bout à l'autre du continent eurasiatique~». Cette vision, également partagée par le ministre des Affaires étrangères Japonais de l'époque, Taro Aso, est alors soutenu par le premier ministre Indien Manmohan Singh puis par le Vice-président Etats-Unien Dick Cheney ce qui aboutit à une première rencontre des membres du Quad en mai 2007 en marge du forum de l'ASEAN à Manille. Ce premier groupement doit cependant rapidement faire face à d'importants défis, certains se posant encore dans sa forme actuelle, qui aboutissent à sa «~dissolution~» moins d'un an plus tard. En effet, l'opposition de la Chine à ce projet est cause d'inquiétude pour plusieurs membres du Quad qui en reste dépendant économiquement, notamment l'Australie dont le premier ministre Kevin Rudd, élu en 2007, refuse de participer au sommet du Quad de janvier 2008. La Corée du Sud, autre allié stratégique des Etats-Unis dans la région, se montre aussi particulièrement sceptique vis-à-vis du Quad qu'elle voit comme trop ouvertement en confrontation avec la Chine et se refuse à le rejoindre.

Cependant en dépit de cet échec, les membres du Quad continuent de renforcer leur coopération via plusieurs accords bilatéraux et trilatéraux~: On peut par exemple citer le Communication Compatibility and Security Agreement signé entre l'Inde et les Etats-Unis en 2015. Sur le plan militaire, les exercices Malabar, lancés en 1992 entre l'Inde et les Etats-Unis, voit la participation du Japon pour la première fois en 2015. De même l'exercice de défense aérien des forces australiennes Pitch Black en 2019 se fait en coopération avec l'armée indienne. Ces élans de coopérations bénéficient par ailleurs d'un contexte favorable avec l'escalade des tensions entre la Chine et l'Inde qui mènent jusqu'à la confrontation armée entre les deux pays au Tripoint de Doklam le long de la LAC à l'été 2017 - le tout appuyé par le retour de Shinzo Abe au poste de premier ministre en 2012. Ainsi le deuxième sommet du Quad à lieu, là encore, à Manille, le 12 novembre 2017. Au programme, l'importance de la dénucléarisation de la Corée du Nord, le maintien de la liberté de navigation en Indo Pacifique et l'émergence du concept de «~free and open~Indo Pacific~» qui est au centre des objectifs du Quad. Ce retour du Quad a, sans surprise, suscité l'ire de la Chine et de la Russie dont le ministre des Affaires étrangères Sergei Lavrov a rejoint les critiques adressées par Wang Yi à l'OTAN et au Quad. La position résolument provocatrice vis-à-vis de la Chine du secrétaire d'État américain sous l'administration Trump, Mike Pompéo, a également contribué à donner du crédit aux revendications chinoises parmi les Etats de la région extérieure au Quad.

Malgré cela, le Quad reste un groupement informel, centré sur une adhésion à des valeurs communes et n'affiche pas de réelles ambitions sur le plan militaire ou institutionnel. Ce «~democratic security diamond~», comme le décrit l'ex premier ministre japonais Shinzo Abe, qui devait s'opposer aux puissances autoritaires en progression dans la région, semble ainsi paradoxalement dépourvu de volet sécuritaire. Le Quad 2021 joint statement, qui précise les principaux objectifs et groupes de travail mis en place dans le cadre du Quad, fait mention d'un groupe dédié à l'effort de vaccination face au Covid 2019, au partage de technologie, et réaffirme la volonté de ses membres à garantir un ordre maritime basé sur le droit international. Néanmoins il n'évoque pas directement la Chine comme étant une menace à cet ordre. Mais il ne s'agit pas pour autant de considérer cet état actuel comme immuable et d'ignorer les différences de visions qui existent entre les membres de cette organisation et qui pourraient amener à une refonte totale de celle-ci en fonction de quelle opinion domine son développement. C'est ce que nous tenterons de clarifier dans l'article suivant en analysant en détail les différentes réalités auxquelles sont confrontés les membres du Quad, quelles conceptions de ce dernier ils défendent et en quoi leurs positions pourraient éventuellement changer face aux crises majeures de ces 5 dernières années.

\textbf{Notes~:}

\textbf{{[}1{]}} Je m'appuie ici sur la définition donnée dans une note de la Maison Blanche sur la stratégie Indopacifique des Etats-Unis publiée en février 2022 et qui rejoint à bien des égards la définition de la «~Broader Asia~» dont parle Shinzo Abe dans son discours sur la Confluence of the two seas. La note en question étant la suivante~: \url{https://www.whitehouse.gov/wp-content/uploads/2022/02/U.S.-Indo-Pacific-Strategy.pdf}

\textbf{{[}2{]}} Le terme utilisé ici est entre guillemets car employé de manière volontairement anachronique. Il provient à l'origine du diplomate américain George F. Kennan en 1946 qui l'a utilisé pour qualifier la stratégie des Etats-Unis vis-à-vis du monde communiste pendant la guerre froide. Il s'agit, en effet, également au cours de cet article d'interroger la pertinence d'une vision bipolaire des rivalités en indopacifique à travers l'éventuelle opposition d'un bloc Sino-russe face à un bloc démocratique qu'incarnerait le Quad.

\textbf{{[}3{]}} citation complète de John Kirby: ~``{[}\ldots{]}And this isn't about an Asian version of NATO.~ NATO is a transatlantic security alliance, the most effective, the most viable one in the world, the most successful one in the world.~». Celle-ci est tirée du document suivant : \url{https://www.whitehouse.gov/briefing-room/statements-releases/2022/06/23/press-briefing-by-press-secretary-karine-jean-pierre-and-national-security-council-coordinator-for-strategic-communications-john-kirby/}

\textbf{{[}4{]}} Il est à noter que la Chine, le Japon et la Corée du Sud sont également membres de l'ASEAN +3, une initiative d'extension de l'ASEAN portée par Singapour, dont la première rencontre s'est tenue en 1997 lors du sommet de Singapour.

\section{Partie 2/3}

Dans l'article précédent, nous avons vu comment le Quad s'est formé et sur la base de quels principes. Cependant la méthode à employer pour accomplir ses objectifs fait l'objet d'importantes oppositions entre les pays membres qui veulent en faire une réelle alliance militaire et stratégique face à la montée de la Chine et ceux qui souhaitent maintenir le Quad dans sa forme actuelle et limiter autant que possible les risques d'affrontement. Ce sont ces courants que nous analyserons ici en considérant le point de vue des différents acteurs engagés, volontairement ou non, dans ce débat qui touche au cœur de ce qui fait la particularité de la zone Indo Pacifique.

Tout d'abord, penchons-nous sur la position des Etats-Unis, principale puissance militaire du Quad et unique en ce qu'elle est relativement éloignée géographiquement des autres membres. Si l'on se réfère à la note de la Maison Blanche {[}1{]} évoquée dans l'article d'avant, ceux-ci ont jusqu'à présent prudemment évité de mentionner le Quad dans leur plan pour renforcer la sûreté de l'espace indopacifique. Effectivement, le Quad n'est mentionné que dans l'axe intitulé «~building connections within and beyond the region~» se concentrant sur son rôle dans la lutte contre les risques épidémiologiques, pour l'aide au développement d'infrastructures et la promotion de l'éducation avec le nouveau Quad Fellowship program{[}2{]}. L'axe «~Bolstering Indo-pacific security~» fait, lui, la part belle aux alliances bilatérales tissées par les Etats-Unis avec les Etats de la région et à la nouvelle alliance trilatérale AUKUS annoncée le 15 septembre 2021 entre les Etats-Unis, l'Australie et le Royaume-Uni. Pourtant, si ce rapport s'ouvre sur une citation de Joe Biden au sommet du Quad tenu seulement quelques jours après l'annonce d'AUKUS c'est bien qu'il occupe une place toute particulière dans la vision qu'ont les Etats-Unis de cet espace. De fait, un sondage réalisé par le Center for Strategic \& International Studies (CSIS), un think tank américain basé à Washington DC, auprès de représentants des élites nationales des pays membres du quad et publié en mars 2020 montre que les Etats Unis sont prêt à soutenir une institutionnalisation du Quad et ne s'opposeraient pas à une vraie politique de coopération militaire. En effet, à la question «~dans quelle mesure seriez-vous prêt à soutenir la création d'un secrétariat permanent du Quad, avec une présidence qui tourne tous les 3 ans~entre les membres ? » les répondants américains se sont exprimés en grande majorité en faveur de cette initiative. Les réponses sont certes plus mitigées face à la question sur la création d'une force armée permanente composée des armées des membres du Quad mais la crise en Ukraine risque fort de faire pencher la balancer en la faveur de cette proposition du côté Etats-Unien. C'est d'autant plus probable que la lutte d'influence entre la Chine et les Etats-Unis dans la région a longtemps été théorisée par les élites américaines et ne cesse de trouver du crédit auprès de ses alliés du «~bloc démocratique~». Pour citer Richard Hass, président du Council on Foreign Relations (CFR) dans la revue Foreign affairs~: «~the rise of authoritarian states challenges the power and leadership of the US~». Une compétition qui entraine aussi les autres états de la région dans sa course comme le souligne Yan Xuetong de l'université de Tsinghua à Pékin, dans le Chinese Journal of International Politics en 2020, rappelant que «~la compétition Sino-américaine crée une pression sur les autres puissances régionales {[}les forçant{]} à prendre parti soit avec les Etats-Unis soit avec la Chine~». Il s'agit toutefois de nuancer cette aspiration à créer un réel bloc démocratique en Asie Pacifique~: Effectivement, bien que la Chine ait multiplié les critiques d'un potentiel «~OTAN asiatique~» en amont du forum de l'OTAN à Madrid en juin 2022, Jens Stoltenberg, secrétaire général de l'OTAN, a rejeté ces accusations en appuyant par une question rhétorique qu'un tel projet n'avait jamais été à l'ordre du jour. Pourtant, si les Etats-Unis disposent encore de nombreuses alliances bilatérales et trilatérales dans la région, la nécessité d'établir une véritable alliance militaire pourrait bien se poser dans un futur proche.

Qu'en est-il pour le Japon à présent~? En tant que pays à l'initiative du Quad, il attache une importance toute particulière à celui-ci et est lui aussi, selon le sondage évoqué précédemment, en faveur d'une institutionnalisation de celui-ci. Toutefois il n'a jamais évoqué le fait d'en faire une alliance militaire et si sa position est certes moins tranchée sur ce sujet que dans le cas de l'Inde, les élites japonaises semblaient écartée une telle évolution à l'époque du sondage. Néanmoins, on peut retrouver un point commun important entre la position japonaise et la position des Etats-Unis qui est l'idée de faire bloc face à la Chine. Dans le cadre du sommet du Quad à Tokyo en mai 2022, le ministre des Affaires étrangères japonais, Hayashi Yoshimasa, a d'ailleurs soutenu la position de plus en plus affirmée des US sur Taiwan. Cela lui a d'ailleurs valu une mise en garde de la part de son homologue chinois suivi par une démonstration de force sino-russe avec le déploiement de 6 bombardiers stratégiques qui ont survolé le territoire japonais. Ces inquiétudes vis-à-vis d'un voisin chinois toujours plus encombrant ont aussi rythmé la visite de Joe Biden à Tokyo en mai de cette année comme le note Jean Yves Colin, membre du conseil d'administration du Think Tank français Asia Center et spécialiste de l'Asie du Nord. A cette occasion, Monsieur Jean Yves Colin tient aussi à rappeler qu'il existe des questions bien plus pressantes pour le Japon que Taiwan, notamment ses inquiétudes en termes de sécurité territoriale concernant les îles Senkaku et l\textquotesingle île Yonaguni. Ainsi, dans ce contexte très tendu, le Japon entend que le Quad joue un rôle majeur pour garantir la souveraineté des Etats sur leur territoire, plus particulièrement dans l'espace maritime Indopacifique. D\textquotesingle ailleurs, lors du sommet quadrilatéral de mai 2022, le Premier ministre japonais Fumio Kishida n\textquotesingle a pas hésité à, selon les termes de Jean Yves Colin, "établir un parallèle entre la situation en Ukraine et l\textquotesingle Indo-Pacifique "{[}3{]}. Cette position semble cohérente avec le concept de «~democratic security diamond~» énoncé par Shinzo Abe en 2006 et reflète bien la vision d'un ensemble fort capable d'imposer ces valeurs dans l'espace indopacifique. Malgré cela, le Japon reste encore particulièrement évasif quant à la nature des moyens à employer pour la réalisation de cet objectif. De fait, la grande majorité de la campagne de coopération militaire entreprise par le Japon se fait dans le cadre de son accord avec les Etats-Unis, et réciproquement. Ainsi le National Security Strategy report publié par le ministère de la défense américain en 2013 dit explicitement que «~le traité de sécurité entre les Etats-Unis et le Japon est la pierre angulaire de la sécurité du Japon~». Un partenariat stratégique que les deux états ont cherché à développer avec, par exemple, la Legislation for Peace and Security en 2015. Il faut également noter que le Japon ne voit pas seulement sa propre sécurité à travers cette alliance mais également celle de ses intérêts dans tout l'espace indopacifique. Shinzo Abe a d'ailleurs qualifié le traité de sécurité EU-Japon de «~pilier indestructible et inébranlable~» qui permet de «~préserver la paix en Asie, en Indo Pacifique et dans le monde entier~»{[}4{]}. Cependant le Japon cherche aussi à diversifier ses partenariats en Asie, notamment l'Inde comme le montre sa participation aux exercices Malabar en 2017. Dès lors, forger une alliance plus vaste et multilatérale pourrait l'aider dans la poursuite de cet objectif.

L'Australie, en revanche, a toujours été très critique vis-à-vis d'un Quad militarisé et plus particulièrement d'un Quad qui ferait bloc face à la Chine. Cela est dû en majeure partie aux forts liens qu'elle entretient avec la Chine dans le domaine économique. En effet, en 2020, la Chine était le principal partenaire commercial du pays, loin devant le Japon et les Etats-Unis, ce qui nuit grandement à son indépendance quand il s'agit de penser la question du Quad. Elle a joué un rôle non négligeable dans l'échec du Quad au début des années 2000 en se retirant de celui-ci en 2008 sous le gouvernement de Kevin Rudd. Néanmoins, de tous les membres du Quad, la position australienne vis-à-vis du Quad est probablement celle qui a le plus évoluée, y compris après son retour en 2017. Comme le décrypte le think tank Japonais, National Institute for Defense Studies(NIDS), l'Australie, qui a toujours pensé l'Océanie comme principale zone d'exercice de sa puissance militaire ainsi qu'économique, a dû faire face à la concurrence de la Chine dans ce même espace et s'est engagé dans une lutte d'influence avec celle-ci. A ce titre, la crise du Covid 19, et plus particulièrement l'aide dispensée aux îles du pacifique par ces deux pays, a été la cause d'importantes tensions entre les 2 Etats. D'un côté, l'Australie a multiplié les initiatives en coopération avec la World Health Organization(WHO) et le World Food Program(WFP), adoptant un modèle résolument multilatéral, de l'autre, la Chine a signé un nombre important d'accord d'aide bilatéraux à destination des pays servant ses intérêts. Cette différence majeure de conception a incité le gouvernement australien à prendre une position bien plus hostile à la Chine que sous Kevin Rudd. En avril 2020, la ministre des Affaires étrangères Australienne Marise Payne a en effet suscité la colère de la Chine en demandant de celle-ci d'avantage de transparence quant à l'origine du virus SARS COV-2 et en demandant la lancée d'une enquête internationale en Chine pour répondre à cette question. La RPC a donc répliqué en imposant une taxation plus forte sur l'orge provenant d'Australie et en mettant en garde les étudiants chinois, principale population international étudiante en Australie avec près de 28\% des effectifs internationaux, allant étudier en Australie qu'il y avait des risques de discriminations. Tout ceci s'est fait dans une atmosphère politique tendue au sein même du gouvernement australien avec la révélation d'affaires de corruption impliquant plusieurs politiciens du gouvernement Rudd et le parti communiste chinois en 2018. Ceci a entrainé le vote d'une loi contre les ingérences étrangères cette même année et a d'autant plus poussé l'Australie à penser son évolution sans la Chine (si ce n'est contre pour certains). Elle a d'ailleurs rejoint les exercices Malabar de 2020 et a renforcé sa coopération économique avec les autres membres du Quad, affichant notamment son souhait de rejoindre le nouveau projet de forum économique en indopacifique IPEF qui regroupera une dizaine de pays de cet espace {[}5{]}. Troy Lee Brown, chercheur en sécurité régionale et maritime à l'université de l'Australie de l'Ouest, a d'ailleurs exprimé que le but pour l'Australie serait à présent «~de proposer une alternative à l'OBOR~» afin de réduire encore d'avantage son attachement à la Chine. Ainsi dans le sondage du CSIS, elle est apparue, au côté des Etats-Unis et du Japon, comme une des puissances motrices du renforcement du quad. Ce changement a également été accompagné d'un changement de ton dans la presse australienne. Là où celle-ci s'est d'abord montré critique eu égard à un Quad fort, certains instituts dont l'Australian Strategic Policy Institute (ASPI) questionnent maintenant la position prudente de l'Inde, que ce dernier qualifie d'ailleurs de «~maillon faible~» du Quad. Il l'explique notamment par sa frontière terrestre importante avec la Chine ce qui la fragilise stratégiquement, et souligne aussi que c'est le seul membre du Quad sans garantie concrète de sécurité provenant des autres membres. Ainsi il estime l'attitude de Delhi comme étant d'avantage trilatérale et comme voyant le Quad comme un simple présentoir plutôt qu'un projet d'envergure face aux ambitions hégémonique chinoise. Un constat qui, semble-t-il, fait échos à la position de l'Australie à l'époque du gouvernement Rudd du point de vue du Japon et des Etats-Unis.

Mais alors qu'en est-il de la position de l'Inde selon les Indiens~? L'avis des élites indiennes recueillie par le CSIS montre en effet que ce pays est de loin le plus hostile à une institutionnalisation du Quad et à une réformation de celui-ci en une alliance stratégique. Cela s'explique pour plusieurs raisons, dont certaines ont déjà été soulevé par l'Australian Strategic Policy Institute~: dans un premier temps, l'Inde souffre d'une dépendance économique majeure avec la Chine et son proche partenaire, la Russie. Une dépendance économique qui se traduit aussi sur le plan militaire avec l'achat massif de matériel russe et chinois par l'armée indienne bien que ces principaux partenaires sur le plan opérationnel soient les Etats-Unis et le Japon avec les exercices Malabar. Dans un même temps, si l'Inde est aussi insistante sur le fait d'éviter la confrontation avec le Chine, c'est qu'elle est déjà sous une forte pression de la part de celle-ci. Les conflits autour de la Line of Actual Control (LAC) ont gagné en intensité depuis le début des années 2000 et ont joué un rôle prépondérant dans la réintégration de l'Inde au Quad en 2017 comme évoqué dans l'article précédent. Cette menace directe à sa souveraineté territoriale est d'ailleurs un des principaux éléments qui distingue l'Inde de l'Australie et l'empêche d'opérer un pivot stratégique aussi significatif que celui réalisé par cette dernière. Ainsi, lors du 4\textsuperscript{ème} sommet Inde-EU à New Delhi en octobre 2020, le secrétaire adjoint américain a fait un discours faisant écho à cette situation bien particulière de l'Inde. Il a ainsi salué la tradition «~d'autonomie stratégique~» qu'a l'Inde, et, s'il a partagé sa vision d'une «~Pax Indo Pacifica~»{[}6{]} en faveur de la protection de la prospérité de l'espace Indopacifique, il a aussi insisté sur le fait que le partenariat de sécurité entre l'Inde et les Etats Unis n'avait pas à impliqué la présence de troupe américaine dans le pays. De même, bien que l'importance de donner à l'Inde les moyens de défendre sa souveraineté ait été au centre de son discours, jamais le Quad n'a été évoqué comme un outil de cette politique de sécurité. On retrouve une idée similaire côté indien dans une prise de parole du Général en chef de l'armée indienne M.M. Naravane dans l'Economic Times en mars 2021 qui a affirmé «~qu'il y aurai définitivement une coopération militaire bilatérale mais que le Quad n'était pas une alliance militaire~».

Est-ce à dire que l'Inde pense uniquement sa défense à travers son alliance bilatérale avec les Etats-Unis et n'a aucun grand projet pour la zone indopacifique~en réponse à la Chine? De toute évidence, l'Inde a su montrer un véritable volontarisme à plusieurs reprises, y compris dans sa façon de penser et de structurer son opposition à la Chine. Sur le plan militaire, ses efforts pour multiplier les exercices militaires avec les autres membres du Quad et travailler son interopérabilité ont, eux, indéniablement joué un rôle de poids dans les progrès du Quad depuis 2017. Sur le plan économique, le projet d'Indo-Japanese Africa Asia Growth corridor (AAGC) porté par le premier ministre indien Narendra Modi en mai 2017 a pour ambition affichée de proposer, à terme, une alternative à l'OBOR. Enfin, l'Inde a été le premier membre du Quad à dénoncer la «~One China policy » menée par Beijing en 2010 en demandant la reconnaissance de la souveraineté de l'Inde sur le Kashmir.

En bref, l'Inde est contrainte d'adopter une position ambiguë dû à sa dépendance vis-à-vis de la Chine mais c'est aussi le pays le plus sensible à l'explosion de la Chine étant donné qu'elle menace directement son intégrité territoriale. Ainsi, plus sa perception de cette menace deviendra importante, plus les chances d'un changement de ton semblable à celui de l'Australie seront grandes, même s'il reste à voir comment celui-ci se fera exactement et si elle placera, ou non, le Quad au centre de ce changement.

En résumé, nous avons vu un peu plus en détail la situation des différents partis impliqués et formulé comment ils envisagent le futur du Quad ainsi que comment ils perçoivent leur propre rôle dans la zone Indo Pacifique. Dans l'article suivant, nous conclurons notre réflexion en posant la question d'une possible extension du Quad et en exposant différent scenarii possibles pour le Quad, y compris celui d'un «~Quad Bloc~» qui aurai un rôle proche de l'OTAN. Nous discuterons, par la même, la probabilité et surtout la légitimité de chacun de ces scenarii afin de mieux en comprendre leur potentiel.

\textbf{Notes~:}

\textbf{{[}1{]}} La note en question~: \url{https://www.whitehouse.gov/wp-content/uploads/2022/02/U.S.-Indo-Pacific-Strategy.pdf}

\textbf{{[}2{]}} Programme lancé le 24 septembre 2021 par le président des États-Unis Joe Biden, conjointement avec les représentants des autres états du Quad, permet aux étudiants provenant de pays membre du Quad (à l'exception des Etats-Unis) et étudiant en master ou en doctorat dans les meilleures universités américaines dans le domaine des sciences, de l'ingénierie, des mathématiques et de la technologie (STEM en anglais) de bénéficier d'une bourse d'étude.

\textbf{{[}3{]}} La citation complète donnée par Jean Yves Colin dans son article est la suivante~: "Russian aggression in Ukraine is a direct threat to the principles of the United Nations. We must not let the same situation happen in the Indo-Pacific"

\textbf{{[}4{]}} Citation extraite d'un discours prononcé par Shinzo Abele 19 janvier 2020 à l'occasion du 60\textsuperscript{ème} anniversaire de la signature du traité de sécurité Japon-EU. Le discours complet est trouvable dans un communiqué du cabinet du premier ministre japonais que voici : \url{https://japan.kantei.go.jp/98_abe/statement/202001/_00002.html}

\textbf{{[}5{]}} Il est intéressant de noter que si la Chine ne fera pas partie de cette initiative, Taïwan ne la rejoindra pas non plus. Ce qui montre bien qu'il y a là une volonté de ne pas risquer la confrontation avec la RPC. Cette position plus ouverte permet l'adhésion de pays réticents à s'opposer ouvertement à la Chine comme la Corée du Sud - cas particulier sur lequel nous reviendront en détail dans le prochain article.

\textbf{{[}6{]}} La citation complète de Stephen E.Biegun est la suivante~: ``One might call this a vision of a Pax Indo-Pacifica, a region at peace, protected and made prosperous in equal measure by those who comprise the Indo-Pacific.~». Le discours dont elle est extraite est à trouver via le lien suivant~: \url{https://2017-2021.state.gov/remarks-by-deputy-secretary-stephen-e-biegun/index.html}

\section{Partie 3/3}

Si nous avons couvert dans l'article précédent les points de vue des différents membres du Quad, il convient à présent de s'interroger sur ce que le Quad pourrait potentiellement devenir à l'avenir et le rôle qu'il pourrait être amené à jouer en Indopacifique comme ailleurs. En lumière de nos analyses précédentes, 2 scenarii bien distincts semblent se détacher~: en premier lieu, l'évolution vers un «~Quad Bloc~», militarisé et surtout institutionnalisé en réponse au bloc créé par la Chine. Puis son opposé, avec un « Quad en retrait~» qui se refuse à afficher la moindre ambition stratégique et reste strictement confiné à un rôle de présentoir pour les démocraties occidentales. Enfin, nous proposerons un troisième scenario, dont la dimension stratégique serait certes avant tout économique, mais qui bénéficierait de réelles prérogatives pour s'imposer en Indopacifique et proposer une alternative concrète au modèle de développement promu par la Chine.

Tout d'abord que serait un «~Quad Bloc~»~? Le terme Bloc fait explicitement référence à des logiques de guerre froide et s'inscrit donc dans la continuité des remarques faites par la Chine à son encontre. Cela suppose un Quad qui énoncerait formellement son objectif de lutter face à cette dernière et se présenterait comme une véritable alliance, plus limitée en taille que l'OTAN, car fondée autour d'un ensemble de pays plus réduit mais aussi plus fortement liés. Mais pour analyser ce scénario plus en détail, nous nous baserons dans un premier temps sur un article de John Hemmings, professeur au Daniel K.Inouye Asia Pacific Center for Security Studies, publié dans le Journal of indo Pacific Affairs en avril 2022, où il argumente en faveur d'une transformation du Quad en une réelle alliance militaire. Ce dernier reprend d'ailleurs une définition particulière du terme «~alliance~», celle de Glenn Snyder {[}1{]} dans Alliance politics, publié en 1997, et qui décrit une alliance comme «~une association formelle d'états pour le recours à la force militaire, dans des circonstances bien déterminées et contre des États qui ne sont pas membres de celle-ci~». Une alliance implique une forme d'engagement de la part des membres et crée une obligation pour ceux-ci de respecter les clauses du traité ou de l'accord secret, ou non, à l'origine de sa formation. Un bon exemple de ce type de formation serait l'alliance entre les Etats-Unis et le Japon ou entre les Etats-Unis et l'Australie. Selon Hemmings, c'est ce type d'alliance qu'exigerai le retour d'une Chine toujours plus ambitieuse et qui inquiète le bloc occidental. Ainsi, en s'appuyant sur les thèses néoréalistes de Kenneth N.Waltz {[}2{]}, l'auteur soutient que les démocraties ont une fâcheuse tendance à ne pas répondre assez fermement pour contrer l'influence de puissances ayant pourtant des prétentions hégémoniques claires. Il établit un parallèle entre la Triple Entente de 1914 et le Quad, soulignant que tous deux sont non-contraignants et rappelle que la Triple Entente n'a pas su répondre dans une mesure suffisante au défi posé par l'avancée allemande à cause de cette hésitation de la part des démocraties. De fait, les preuves de l'ambition de la Chine se multiplient comme le notent plusieurs rapports stratégiques émis par les gouvernements des membres du Quad. Par exemple en 2018, le rapport US National Defense Strategy parlait d'une RPC qui «~poursuit son programme de modernisation militaire et qui cherche à atteindre une hégémonie régionale en indopacifique~». Le Defense White Paper publié par le ministère de la défense japonais en 2021 et la Defense Strategic Update publiée par son équivalent australien un an plus tôt font d'ailleurs des constats similaires. De son côté, Xi Jinping a exhorté à plusieurs reprise - en mars 2021, janvier 2021 ou encore mai 2020 - la RPC à se préparer à la guerre à tout moment.

Le développement d'un Bloc Sino-Russe est aussi un facteur important à prendre en compte pour comprendre l'équation complexe posée par la coopération militaire en indopacifique. Pour Andrew Yeo, Senior fellow et titulaire de la chaire de la Fondation SK-Korea en études coréennes à Brookings, l'invitation de l'OTAN aux pays de l'Indo pacifique s'inscrit dans cette logique et il estime qu'elle «~témoigne de la prise de conscience d'un lien plus profond entre la sécurité {[}\ldots{]} transpacifique, {[}\ldots{]} en raison de la coopération accrue entre la Chine et la Russie~». Un rapprochement qui prend place depuis de nombreuses années sur le plan économique comme militaire avec par exemple la participation de la Chine à l'exercice militaire annuel russe depuis 2018. Il s'accompagne d'un processus d'éloignement du bloc occidental qui n'a lui aussi cessé de s'accélérer et atteint aujourd'hui des sommets avec la guerre en Ukraine {[}3{]}.

Malgré tout cela, ils existent aussi de puissants contre arguments à cette vision. Premièrement, la viabilité d'une réelle alliance sino-russe est loin d'être évidente. Si la Russie n'a pour l'instant pas d'autres choix que de se tourner vers la Chine, elle se retrouve dans une position particulièrement complexe, étant incapable de négocier d'égal à égal face au géant chinois. Jae Seung Lee, professeur et titulaire de la Chaire Jean Monnet à l'université de Corée, n'a d'ailleurs pas manqué de soulever, à l'occasion d'une conférence tenue à Paris à maison Suger, qu'il serait particulièrement compliqué pour la Russie de faire valoir son point de vue dans un monde dominé par la Chine et que sa position de fragilité risque d'être rapidement exploitée par cette dernière. Dès lors, même si cela semble impossible pour l'instant, si la Russie veut retrouver son indépendance et peser dans le monde elle devra se trouver de nouveaux partenaires, y compris dans la zone Indopacifique. Deuxièmement, comme explicité dans l'article précédent, la perspective d'un «~Quad Bloc~» est très loin de faire l'unanimité parmi ses membres, pour ne pas dire qu'elle est purement et simplement inenvisageable à l'heure actuelle. A moins qu'un évènement majeur ne survienne et amène les plus réticents à réévaluer le seuil de menace présenté par la Chine et à y apporter une réponse militaire, il est fort peu probable qu'une militarisation du Quad s'opère au cours de la décennie à venir. Même si la crise en Ukraine a été un catalyseur pour la remilitarisation de l'Europe et du Pacifique, elle ne suffira vraisemblablement pas à pousser les États plus loin que leurs ambitions d'alliances bilatérales ou trilatérales. Enfin, faire du Quad un acteur fédérateur du bloc démocratique en indopacifique serait, dans ce cas de figure, éminemment délicat. Un bon nombre de pays en Asie Pacifique sont encore très dépendants de la Chine économiquement et seront forcément beaucoup plus hésitants à rejoindre une alliance ouvertement en conflit avec celle-ci. Cela pourrait aussi poser d'autres défis stratégiques aux répercussions dépassant l'espace indopacifique. A ce titre, le cas de la Corée du Sud est particulièrement intéressant. Jae Seung Lee a, au cours de la même conférence, insisté sur le fait que la Corée du Sud ne chercherait jamais à rejoindre une alliance visant la confrontation avec la Chine pour plusieurs raisons~: La première, d'ordre économique, est lié à des enjeux de supply chain -- la Chine fournissant l'écrasante majorité des terres rares nécessaires à la production de semi-conducteurs et autres produits de haute technologie dans laquelle la Corée du Sud excelle. La deuxième, d'ordre géopolitique, est que pour négocier avec une Corée du Nord toujours plus agitée, et disposant d'un programme nucléaire en progrès constant {[}4{]}, il est nécessaire de préserver de bonnes relations avec la Chine dont la Corée du Nord dépend complètement. La dénucléarisation de la Corée du Nord, d'ailleurs évoquée dans le sommet du Quad à Tokyo cette année, semble donc irréalisable dans le cadre d'un «~Quad Bloc~». En résumé, un tel développement apparait comme très peu probable à l'heure actuelle. Il est en revanche probable qu'il ne trouve pas de réelles justifications dans les années à venir, sous réserve d'évènements exceptionnels.

Est-ce à dire pour autant qu'un «~Quad en retrait~» comme évoqué en introduction lui serait préférable~? Une vision à l'opposé exact du «~Quad Bloc~» reviendra indéniablement à un retour en arrière pour le Quad ou du moins à une forme de continuité avec ce qu'il est actuellement. Or dans son état actuel, ses prérogatives sont pour le moins limitées voire inexistantes étant donné qu'il est un groupe informel et non une institution à part entière. Cela limite grandement son champ d'action et si, dans le domaine économique, il a déjà mis en place plusieurs programmes d'aide qui sont déjà effectifs, ces initiatives font encore pâle figure par rapport à des associations comme l'ASEAN dont le pouvoir, y compris dans ce domaine, reste pourtant tout à fait relatif. C'est encore plus vrai dans le domaine militaire où la quasi-intégralité des avancées se font au niveau bilatéral ou trilatéral ce qui donne l'impression que l'influence du Quad n'est guère que celle du plus puissant de ses membres et n'apporte donc aucune réelle valeur ajoutée. Le 3 juin 2022, le Quad a lancé l'Indo-Pacific partnership for Maritime Domain Awareness (IPMDA) qui prévoit le déploiement de satellites et d'autres système de surveillance pour lutter contre la pêche illégale -- dont la Chine a fait une spécialité en mer de Chine Septentrionale et Méridionale -- ainsi que le partage d'information avec des pays partenaires tels que Singapour ou les îles Salomon. Une initiative d'apparence ambitieuse que la Chine, par l'intermédiaire de son ministre des affaires étrangères, n'a d'ailleurs pas manqué de critiquer en soutenant que le but du Quad avec l'IPMDA était de «~contenir la Chine~» et de «~faire des autres pays d'Asie Pacifique des pions de l'hégémonie américaine~». Cependant cela reste une mesure essentiellement préventive et la capacité des états engagés dans cet effort à lutter concrètement contre les incursions répétées dans leur ZEE reste à prouver. Dès lors, si ce type de Quad trouverait sans aucun doute plus facilement des candidats pour participer à son expansion, il jouerait principalement le rôle de vitrine des démocraties plutôt que celui d'un acteur capable de s'imposer en indopacifique et serait donc incapable de répondre aux grands défis auxquels la région fait face.

Qu'elle pourrait alors être une voie alternative à ces deux extrêmes~? Certes, faire du Quad une véritable alliance militaire semble compliqué et peu souhaitable, mais il n'en garde pas moins un potentiel certain en termes de stratégie économique. A condition de limiter le nombre d'États membres, il pourrait permettre une intégration économique bien plus forte des états principaux de la zone et les aider à se détacher, au moins en partie, de leur dépendance économique avec la Chine. Une telle évolution serait loin d'être inutile sur le plan géopolitique en offrant davantage de latitude à des pays comme l'Inde pour penser leur positionnement y compris sur le plan militaire. De même, un effort d'intégration économique et technologique ambitieux de la part du Quad pourrait fournir l'élan nécessaire à d'autres formes de coopération économique en plein développement, à l'image de l'IPEF ou de l'AAGC, pour concurrencer la «~One Belt One Road~» initiative de la RPC. Pour revenir à la Corée du Sud, si une dynamique d'opposition plus frontale à la Chine pourrait la refroidir dans l'idée de rejoindre un «~Quad +~» cela pourrait aussi lui permettre de faire face plus sereinement à ses problèmes de supply chain. Selon le CSIS, faire partie du Quad lui permettrait aussi de faire valoir ses positions sur la scène régionale et internationale, y compris concernant la Corée du Nord, en s'appuyant sur les Etats-Unis. Cependant le CSIS reconnait aussi que cette intégration est nécessairement conditionnée à la garantie d'un soutient indéfectible de la part de ces derniers étant donné les risques d'une réaction chinoise. La Chine ne donnera certes pas automatiquement une réponse concrète sur le plan économique ou militaire, mais si une telle chose devait se produire il est impératif pour la Corée du Sud et pour tout autre pays aspirant à rejoindre le Quad d'avoir la protection de Washington et des autres membres du Quad. Si l'administration Trump et ses nombreuses inconsistances en termes de politique étrangères ont durablement érodé la confiance placée en les Etats-Unis par ses alliés, dont la Corée du Sud, la promotion par Biden d'un «~free, open and rules-based Indopacific~» semble plus en accord avec cette nécessité. Toutefois, comme le souligne le journal le Monde dans un article publié le 29 juin 2022, «~l'importance du Quad pour Biden~est encore difficile à cerner » et la valeur qu'il lui accordera face à d'autres projets concurrents comme le partenariat AUKUS doit être clarifié afin de pouvoir penser sereinement l'avenir du Quad.

En conclusion, la forme que prendra le Quad à l'avenir est subordonnée à bien des facteurs extérieurs au contrôle des États, mais ce sont leurs prises de position à l'échelle globale qui détermineront en dernier lieu la direction que prendra son évolution. Sans un soutien clair et un authentique volontarisme des Etats-Unis, il est peu probable qu'il puisse s'institutionaliser et encore moins devenir un acteur significatif de l'ordre géopolitique indopacifique. De même, si l'Inde et les pays susceptibles de rejoindre le Quad n'estiment pas la menace présentée par la Chine comme supérieure aux bénéfices liés à la préservation de leur intérêt économique avec celle-ci, il est radicalement impossible que le Quad prenne bien plus d'envergure qu'il n'en a actuellement. Cela étant dit, bien que la probabilité d\textquotesingle un OTAN asiatique basée sur le Quad soit effectivement faible, si les pays de l\textquotesingle Asie-Pacifique sont disposés à faire confiance à un Quad plus fort et mieux organisé, ils pourraient bien en bénéficier.

\textbf{Notes~:}

\textbf{{[}1{]}} La citation originale de Glenn Snyder, professeur émérite en science politique à l'université de Caroline du Sud, pour définir une alliance est la suivante~: ``formal associations of states for the use of military force, in specified circumstance against states outside their own membership''

\textbf{{[}2{]}} Le concept de néoréalisme en géopolitique a été développé par Kenneth N. Waltz, un des politologues les plus reconnus dans le domaine des relations internationales et membre de la faculté de l'université de Berkeley et de Columbia, dans son ouvrage \emph{Theory of international politics} publié en 1979. Cette théorie laisse peu de place à la coopération dans les relations internationales et met l'accent sur la compétition et le conflit comme des composantes majeurs de l'ordre géopolitique.

\textbf{{[}3{]}} Il est intéressant de noter que déjà en 2021 lors d'une visite en Chine du ministre des affaires étrangères russe, Sergeï Lavrov, ce dernier a exprimé l'importance de sortir du système Swift en soulignant sa volonté de «~renforcer l'indépendance technologique de la Russie et de s'écarter du système de paiement international contrôlé par l'Occident.~»

\textbf{{[}4{]}} Concernant les avancées majeures réalisées par la Corée du Nord, ces dernières années, dans le domaine du nucléaire militaire, on peut citer~: le développement de l'IRBM (intermediate-range ballistic missile) Hwasong-12 en janvier 2021, capable de toucher l'Alaska, ou encore , plus récemment, de l'ICBM (intercontinental ballistic missile) Hwasong-15 qui pourrait potentiellement frapper la métropole Etats-Unienne.

\section{Sources pour les trois articles}

\begin{itemize}
  \item \textbf{Japanese Press}
  \begin{itemize}
    \item \url{https://english.kyodonews.net/news/2022/06/1345330ff0f8-no-asian-nato-says-us-as-alliance-brings-japan-s-korea-to-summit.html}
    \item \url{https://www.japantimes.co.jp/news/2022/06/20/national/japan-korea-australia-new-zealand-nato-meeting/}
    \item \url{https://www.japantimes.co.jp/news/2022/06/29/national/japan-sk-us-talks/}
    \item \url{https://asia.nikkei.com/Politics/International-relations/Japan-will-turn-to-Quad-in-new-Cold-War-Defense-Ministry-think-tank}
  \end{itemize}
  \item \textbf{British Press}
  \begin{itemize}
    \item \url{https://www.bbc.com/news/world-asia-india-61547082}
  \end{itemize}
  \item \textbf{North American Press}
  \begin{itemize}
    \item \url{https://thediplomat.com/2017/11/why-has-australia-shifted-back-to-the-quad/}
    \item \url{https://foreignpolicy.com/2020/10/08/why-quad-alliance-trump-cares-about-india-australia-japan-china/}
    \item \url{https://www.voanews.com/a/japan-s-kishida-considers-joining-nato-summit--sources/6603976.html}
    \item \url{https://www.voanews.com/a/what-to-expect-from-south-korea-s-first-nato-summit/6635935.html}
    \item \url{https://edition.cnn.com/2021/03/11/asia/quad-us-india-japan-australia-intl-hnk/index.html}
    \item \url{https://time.com/6188414/japan-south-korea-nato-summit/}
  \end{itemize}
  \item \textbf{Australian Press}
  \begin{itemize}
    \item \url{https://www.aspistrategist.org.au/india-and-the-quad-weak-link-or-keystone/}
    \item \url{https://www.news.com.au/technology/innovation/the-quad-launches-bold-plan-to-combat-chinas-dark-ships/news-story/4aedcc8c4403199b8bd5ab69c410305a}
  \end{itemize}
  \item \textbf{French Press}
  \begin{itemize}
    \item \url{https://www.lemonde.fr/en/international/article/2022/06/29/japan-and-south-korea-want-nato-to-look-toward-asia_5988383_4.html}
    \item \url{https://www.france24.com/fr/asie-pacifique/20210924-quad-aukus-comment-washington-isole-la-chine-dans-la-zone-indo-pacifique}
  \end{itemize}
  \item \textbf{Indian Press}
  \begin{itemize}
    \item \url{https://www.thehindubusinessline.com/opinion/can-the-quad-rise-to-be-an-asian-nato/article32103609.ece}
    \item \url{https://www.hindustantimes.com/india-news/top-us-diplomat-calls-for-expansion-of-quad-group/story-R06qlCO6ft59K6GMJPkEPM.html}
    \item \url{https://www.ndtv.com/india-news/pm-narendra-modis-at-quad-summit-in-tokyo-quad-a-force-for-good-helps-make-a-better-indo-pacific-region-3003210}
    \item \url{https://economictimes.indiatimes.com/news/defence/indias-fine-balancing-act-with-quad-and-brics-meet-in-new-york/articleshow/71338616.cms}
    \item \url{https://www.ndtv.com/world-news/quad-meet-quad-meeting-in-tokyo-checking-china-with-asian-nato-what-is-the-quad-alliance-3003474}
  \end{itemize}
  \item \textbf{Chinese Press}
  \begin{itemize}
    \item \url{https://www.scmp.com/week-asia/explained/article/3122933/what-quad-and-how-will-it-impact-us-china-relations-under-biden}
  \end{itemize}
  \item \textbf{Reports and briefing from think tanks, specialised institutes and intergovernmental institutions}
  \begin{itemize}
    \item \url{https://www.cfr.org/in-brief/quad-getting-more-ambitious-indo-pacific}
    \item \url{https://csis-website-prod.s3.amazonaws.com/s3fs-public/publication/200312_BuchanRimland_QuadReport_v2\%5B6\%5D.pdf?fuRA6mwjWYKqROtSmJD4u5ct.vijdkZv}
    \item \url{https://www.airuniversity.af.edu/JIPA/Display/Article/2979986/should-the-quad-become-a-formal-alliance/}
    \item \url{http://www.nids.mod.go.jp/english/publication/east-asian/pdf/2021/e01.pdf}
    \item \url{http://www.nids.mod.go.jp/english/publication/east-asian/pdf/2021/e07.pdf}
    \item \url{https://www.csis.org/analysis/return-quad-will-russia-and-china-form-their-own-bloc}
    \item \url{https://www.csis.org/analysis/south-korea-would-benefit-joining-quad}
    \item \url{https://asiacentre.eu/fr/2022/05/25/le-retour-du-pivot-nord-asiatique-visite-de-j-biden-a-seoul-et-tokyo/}
    \item \url{https://www.indembassy-tokyo.gov.in/public_files/assets/pdf/Quad_Summit_Fact_Sheet_english.pdf}
  \end{itemize}
  \item \textbf{Reports from governmental organisations}
  \begin{itemize}
    \item \url{https://www.whitehouse.gov/wp-content/uploads/2022/02/U.S.-Indo-Pacific-Strategy.pdf}
  \end{itemize}
  \item \textbf{Official speeches}
  \begin{itemize}
    \item \url{https://www.mofa.go.jp/region/asia-paci/pmv0708/speech-2.html}
    \item \url{https://www.nato.int/cps/en/natohq/opinions_197292.htm}
  \end{itemize}
  \item \textbf{Conference}
  \begin{itemize}
    \item Conference The Challenge of Eurasian Geopolitics and Korea’s New Partnership with Europe, 24 June 2022, held by Jae Seung Lee, Professor and Jean Monnet Chair at Korea University and moderated by Jean-Yves Colin, researcher at the Asia Centre and expert on North East Asia.
  \end{itemize}
  \item \textbf{Academic works}
  \begin{itemize}
    \item Rémi Hémez, \emph{L’alliance militaire entre la Corée du Sud et les Etats-Unis Sous Moon Jae-In, 2017-2022}, Revue défense nationale, n°850 mai 2022
    \item Kenneth N. Waltz, \emph{Theory of international politics}, New York, Mac Graw Hill, 1979
    \item Glenn Snyder, \emph{Alliance politics}, Ithica: Cornell university press, 1997
  \end{itemize}
\end{itemize}

\end{document}
